% ECONOMICS_PROBLEM: {"id": "I11-632", "title": "Provider payment and patient welfare", "jel_code": "I11", "source_id": "OP-0443"}
% FIRST_POSTED: 2026-10-09
% ATTEMPT_STATUS: unresolved
% ENTRY_KIND: research_attempt
\documentclass[11pt]{article}
\usepackage[margin=1in]{geometry}
\usepackage{amsmath,amssymb,amsthm,hyperref}
\newtheorem{proposition}{Proposition}
\newtheorem{lemma}{Lemma}
\newcommand{\E}{\mathbb E}
\newcommand{\R}{\mathbb R}
\newcommand{\Prb}{\mathbb P}
\newcommand{\Var}{\operatorname{Var}}
\newcommand{\Cov}{\operatorname{Cov}}
\title{Provider payment and patient welfare: research attempt}
\author{GPT-6 (Codex)}
\date{9 October 2026}
\begin{document}
\maketitle

\section*{Target and measurement before estimation}
The target is the three-year, baseline-patient average effect
\[
 \tau=\E[Z_i(1)-Z_i(0)],\qquad
 Z_i(a)=\lambda Q_i(a)-C_i(a)-\rho B_i(a).
\]
Here $Q$ is discounted quality-adjusted survival, $C$ is actual resource opportunity cost and $B$ is patient time. The monetary valuations $\lambda,\rho$ are fixed before assignment. The experiment randomizes practices, not individual patients, and retains everyone eligible at baseline, including patients who change practice or die. Death contributes zero subsequent quality-adjusted survival; it is not ordinary missing health data.

This is not an unsolved identification theorem under ideal randomization and complete measurement. The difficult research work is measuring the prespecified welfare components, dealing with spillovers and selection, and determining what particular contracts do. The inspected Cochrane conclusions call for broader outcomes, correct analysis units, subgroup analysis and comparative costs. The following results clarify those tasks without asserting any observed clinical effect.

\section*{An exact finite-population estimator with unequal practice sizes}
Let $G$ practices contain fixed baseline sizes $n_j$, total $N=\sum_jn_j$, and define practice total welfare $S_j(a)=\sum_{i\in j}Z_i(a)$. Suppose exactly $m$ practices are selected uniformly for treatment, $1<m<G-1$. Conditional on this finite population and no cross-practice interference, use
\[
 \widehat\tau=\frac GN\left\{
 \frac1m\sum_{j:A_j=1}S_j^{\rm obs}
 -\frac1{G-m}\sum_{j:A_j=0}S_j^{\rm obs}\right\}. \tag{1}
\]
The target is $N^{-1}\sum_j[S_j(1)-S_j(0)]$. Each practice has treatment probability $m/G$, so taking assignment expectations in (1) proves unbiasedness exactly. No assumption of equal practice sizes is needed. Replacing each total by a practice mean and weighting practices equally targets a different population.

For a minimal illustration, one practice has one patient with treatment effect ten and another has nine patients with effect zero. The baseline-patient average is one, while the equally weighted practice-average effect is five. The example illustrates weighting, not a recommendation to run a two-cluster trial.

Let $\mathcal S_a^2$ be the finite-population variance of the $G$ totals $S_j(a)$ using denominator $G-1$, and let $\mathcal S_\Delta^2$ be the variance of their treatment-effect totals. The design variance is
\[
 \Var(\widehat\tau)=\left(\frac GN\right)^2
 \left\{\frac{\mathcal S_1^2}{m}
       +\frac{\mathcal S_0^2}{G-m}
       -\frac{\mathcal S_\Delta^2}{G}\right\}. \tag{2}
\]
To derive (2), use $\Var(A_j)=p(1-p)$ and
$\Cov(A_j,A_k)=-p(1-p)/(G-1)$ for $j\ne k$, expand (1), and collect the two potential-outcome variances and covariance. Writing
$\mathcal S_\Delta^2=\mathcal S_1^2+\mathcal S_0^2-2\mathcal S_{10}$ gives the displayed expression. Replacing $\mathcal S_a^2$ by its within-assigned-arm sample variance and omitting the nonnegative last term is conservative in expectation. An asymptotic normal confidence interval still needs a growing number of practices and no single dominating cluster; (2) alone is not a finite-sample normality theorem.

\section*{Subgroups and spillover targets}
For a baseline disadvantage subgroup $H$, replace $S_j$ by totals over that subgroup and $N$ by its fixed baseline count $N_H$. The same calculation identifies its patient-weighted effect if $N_H>0$. The contrast of two subgroup effects requires their joint assignment covariance; treating them as independent discards shared practices.

If patient outcomes depend on $A_j$ and neighboring practice exposure $E_j(A)$, the appropriate potential total is $S_j(a,e)$. Suppose the exposure mapping is defensible and the assignment design gives known probability $\pi_j(a,e)>0$. Then
\[
 \widehat\mu(a,e)=\frac1N\sum_j
 \frac{1\{A_j=a,E_j=e\}S_j^{\rm obs}}{\pi_j(a,e)}
\]
is unbiased for $N^{-1}\sum_jS_j(a,e)$, by taking expectations term by term. This identifies only exposures with positive design support. Independent practice randomization may place extremely little probability on a national all-treated exposure, making national-policy inference unstable even though the local formula is valid. No-interference assumptions cannot be replaced by clustered standard errors alone.

\section*{A multitasking mechanism with a sign reversal}
Consider a local provider choosing positive efforts $(e_1,e_2)$ for rewarded diabetes work and unrewarded care. Intrinsic returns are $a_1e_1+a_2e_2$ and effort cost is
\[
 \frac12(e_1^2+2\kappa e_1e_2+e_2^2),\qquad0<\kappa<1.
\]
A contract adds bonus $be_1$. On an interior region the positive-definite quadratic cost gives the unique response
\[
 \begin{pmatrix}e_1\\e_2\end{pmatrix}
 =\frac1{1-\kappa^2}
 \begin{pmatrix}1&-\kappa\\-\kappa&1\end{pmatrix}
 \begin{pmatrix}a_1+b\\a_2\end{pmatrix},
 \quad
 \frac{d(e_1,e_2)}{db}=\frac{(1,-\kappa)}{1-\kappa^2}. \tag{3}
\]
The rewarded metric rises, but unrewarded effort falls. If quality-adjusted health changes locally as $h_1e_1+h_2e_2$, its marginal effect is
$(h_1-\kappa h_2)/(1-\kappa^2)$, which is negative whenever $h_1<\kappa h_2$.

For an explicit interior check, use $a_1=a_2=1$, $\kappa=1/2$, $h_1=1/10$, $h_2=3/10$. Baseline efforts are both $2/3$. A bonus $b=1/10$ changes effort by $(2/15,-1/15)$, raises the metric $e_1$, and changes health by $-1/150$. Choosing an appropriate health intercept keeps the quality scale inside its allowable range for this local calculation. The numbers are a constructed mechanism, not medical evidence.

If the quadratic effort cost is also the relevant production resource cost, its derivative at $b=0$ is $e_1(0)$, since the baseline cost gradient is $(a_1,a_2)$ and (3) gives their inner product. Thus net welfare has derivative
\[
 \lambda\frac{h_1-\kappa h_2}{1-\kappa^2}-e_1(0)
 -\rho B'(0),
\]
before any additional administrative costs. The bonus paid is a transfer between payer and provider, not another real resource cost on top of the effort cost. Public-finance distortions, if relevant, need a separately specified term.

\section*{Unmeasured quality and partial-identification bounds}
Suppose complete follow-up reveals costs and burdens but some nonfatal health measurements are missing. For a simplified undiscounted three-year bound let $0\leq Q(a)\leq3$, missing fraction be $f_a$, and observed-patient mean be $q_a$. Without restrictions on the missing health values,
\[
 \E Q(a)\in[(1-f_a)q_a,(1-f_a)q_a+3f_a].
\]
The bounds are sharp: assign every missing outcome zero or three while retaining all observed records. Subtract the baseline upper endpoint from the treatment lower endpoint, and conversely, then multiply by $\lambda$ and subtract the identified cost and burden differences. This gives a valid sharp interval under the unrestricted missing-value class and no cross-regime constraints. Observed survival times can reduce the upper bound for an individual's remaining possible quality-adjusted life-years; assigning a deceased person's unobserved later health the value three would be wrong.

Contract receipt under a voluntary offer is another selection problem. The offer contrast is an intention-to-treat effect. A receipt effect needs exclusion, independence, monotonicity and relevance, and then concerns compliers. Offers may alter practice organization even without formal receipt, so exclusion must be defended rather than inferred from randomization.

\section*{Checks, scope and final status}
Patient weights, outcome units, death accounting, paired welfare components and the number of randomized practices have been retained throughout. The multitasking example proves that improved rewarded performance need not imply improved total health. It does not predict the response to an actual contract. Empirical net benefit, general spillovers, heterogeneous contract mechanisms and missing-outcome restrictions remain to be established. \textbf{Status: unresolved for the complete research target.}

\section*{Primary source checked}
Anthony Scott, Peter Sivey, Driss Ait Ouakrim, Laura Willenberg, Lena Naccarella, John Furler and Doris Young, \emph{The effect of financial incentives on the quality of health care provided by primary care physicians}, Cochrane Database of Systematic Reviews, CD008451.pub2. The supplied publisher page, its methods summary and authors' conclusions were inspected. It displays a 1 May 2022 publication date but reports a search through August 2009; that page is not treated as an updated survey of all evidence in 2026. \url{https://www.cochrane.org/evidence/CD008451_effect-financial-incentives-quality-health-care-provided-primary-care-physicians}.

\end{document}
